\documentclass[11pt,a4paper]{article}
\usepackage[T1]{fontenc}
\usepackage{lmodern,amsmath,amssymb,booktabs,longtable,array,geometry,hyperref}
\geometry{margin=23mm}
\hypersetup{hidelinks}
\emergencystretch=3em
\newcommand{\C}{\textbf{(c)}\ }
\newcommand{\D}{\textbf{(d)}\ }
\newcommand{\dd}{\mathrm d}
\newcommand{\sech}{\operatorname{sech}}
\newcommand{\file}[1]{\nolinkurl{#1}}
\newcommand{\alg}{\file{shiftcharge_algebra.py}}
\newcommand{\modes}{\file{shiftcharge_modes.py}}
\newcommand{\calcheck}{\file{shiftcharge_calibration.py}}
\newcommand{\proj}{\file{shiftcharge_projection.py}}
\newcolumntype{L}[1]{>{\raggedright\arraybackslash}p{#1}}
\title{One-loop shift charge of Boussinesq solitons:\
centroid identity, mean-profile response, and the unresolved two-body coefficient}
\author{Quantum Boussinesq solitons project --- PI: Gustav Delius}
\date{4 November 2026}
\begin{document}
\maketitle
\begin{abstract}
\C The exhibited freedom $c(P_1M_2-P_2M_1)$ has one coefficient and admits
a rank-zero defect consistency test, but the cited package does not classify
the entire continuum ambiguity. The centroid identity follows from the
canonical continuity equation, subject to its domain and boundary limits.
It does not justify holding the kink height fixed while moving along the
quantum one-soliton dispersion branch.

\C A fresh continuum-mode calculation gives a finite conditional mean-profile
shift charge, including the finite quadratic counterterm. Its direct momentum
insertion agrees with the canonical derivative of the audited one-body trace.
At fixed momentum the corresponding kink-height correction is nonzero;
its contribution cancels the apparent leading rational term from changing
the shift charge alone. The separate, already-known Darboux recoil reproduces
Toda's $2(r-r^{-1})/\pi$, but is not a constant times $P\wedge M$.
\D \textbf{Conclusion (C):} neither $c_{\rm bad}$ nor
$c_{\rm good}(\lambda)$ is determined. The proposed identification fails as
a complete mechanism; the boundary-compatible projected expectation,
two-body determinant and common quantum frame remain missing.
\end{abstract}

\section{Scope, evidence and the freedom actually established}
Labels mean literature proof (a), literature heuristic (b), project derivation
or computation (c), and an unestablished physical identification (d).
No external full text or abstract was newly read and no external reference
identifier is asserted. Internal sources are listed at the end. Every new
displayed (c) formula is covered by the named scripts, located in
\file{phase2/scripts/}; their accompanying JSON files retain the numerical
results and explicit nulls. Conditional mathematical results are not promoted
to constructed quantum observables.

\C The fixed good convention is retained, with the bad equation used only
for formal continuation and Toda validation. README Status, the master
report including its ranked programme, the named packages, their relevant
audits and the errata were consulted. The frame and determinant qualifications
override package interpretations. Neither \file{audit_defect_bootstrap.tex}
nor \file{audit_good_renorm.tex} existed when checked during this work;
a reference to the latter in another audit is not evidence of its contents.
No old script was executed or imported; no git operation or existing-file edit
was made. No renormalisation condition is selected for the PI.

\paragraph{Independent structure check \C (\alg).}
Write $Q_{12}=P_1M_2-P_2M_1$. In the displayed two-function ansatz,
\begin{equation}
 F_{12}=A P_1M_2+B M_1P_2,\qquad F_{12}+F_{21}=(A+B)(P_1M_2+P_2M_1).
\end{equation}
Thus antisymmetry leaves one coefficient, $A=-B=c$. All mixed $3\times3$
minors vanish, while mixed $2\times2$ minors need not vanish. This is a
rank-two example, not an exhaustive classification. In particular,
\file{defect_bootstrap.tex} explicitly retains additional unknown $U_{12}$.
The stronger premise ``the full phase is known up to one number'' is not
established by that package or the audits.

\C For a good KP chord with endpoints $a,b$,
\begin{gather}
 k=a-b,\quad v=\sqrt3(a+b),\quad a^2+ab+b^2=\tfrac14,\qquad
 M=m(a)-m(b),\quad P=p(a)-p(b),\label{chords}\\
 m(z)=2z,\qquad p(z)=2\sqrt3 z^4-z^2/\sqrt3-\sqrt3/36.
\end{gather}
Polynomial reduction verifies the charge primitives independently. They
telescope on a triangle. For a spectator $t$, the lift
$W_z(t)=c[P_t m(z)-M_t p(z)]$ satisfies
$W_b-W_a=-F(t,a\!\to\! b)$, with the orientation reversed if the chord is
reversed. This supplies a rank-zero witness for the homogeneous cycle
constraint; scalar transmission commutativity adds no rank. It does not
construct a quantum absorption vertex. The corrected Toda characteristic
source and its axis data, discussed in \file{twobody_toda.tex} and its audit,
are additional dynamical inputs, not continuum bootstrap consequences.

\section{What changes at fixed physical momentum}
\paragraph{Conventions and inversion \C (\alg).}
For the good branch put
\begin{gather}
 k=\sin\alpha,\quad v=\epsilon\sqrt{1-k^2},\quad
 M_0=2k,\quad \Pi_0=-2kv,\quad P_0=\tfrac23vk^3,\\
 E_0=\tfrac2{15}k^3(5-4k^2),\qquad \dd E_0=v\dd P_0.
\end{gather}
Take a label $a$ away from $P_0'(a)=0$ and write
$P=P_0+\hbar p$, $\Pi=\Pi_0+\hbar q$, $M=M_0+\hbar m$.
At fixed physical $P$, the argument changes by $-\hbar p/P_0'$; hence
\begin{equation}
 \bar q=q-\frac{\Pi_0'}{P_0'}p,\qquad
 \bar m=m-\frac{M_0'}{P_0'}p.\label{fixedP}
\end{equation}
The bars are essential: fixed classical label, fixed velocity, fixed kink
height and fixed physical momentum are different operations.

\paragraph{Topological term and frame subtraction \C (\alg).}
In the specified chamber-reversing convention the tree phase contains
\begin{equation}
 T_0=-\tfrac23(\Pi_{0,1}M_{0,2}-\Pi_{0,2}M_{0,1}),\qquad
 S_{\rm cb}=\tfrac12(\Pi_1M_2+\Pi_2M_1),\quad
 \delta^{\rm raw}=\delta+S_{\rm cb}.
\end{equation}
Replacing its charge arguments gives the following \emph{part} of the
order-one exponent:
\begin{align}
 T_1|_P&=-\tfrac23[\bar q_1M_{0,2}+\Pi_{0,1}\bar m_2
                  -\bar q_2M_{0,1}-\Pi_{0,2}\bar m_1],\label{T1}\\
 (S_{\rm cb})_1|_P&=\tfrac12[\bar q_1M_{0,2}+\Pi_{0,1}\bar m_2
                  +\bar q_2M_{0,1}+\Pi_{0,2}\bar m_1].\label{S1}
\end{align}
The latter must be subtracted from a raw fixed-chamber answer; it is symmetric
and must not be counted again as an antisymmetric phase.
\D Substitution into $T_0$ alone is not a derivation of the full quantum phase:
the dynamical term, quantum delays and state map can also change.

\C If $\bar m=0$, a sufficient condition for $T_1=cQ_{12}$ is
\begin{equation}
 \bar q(P)=\kappa P+aM_0(P),\qquad c=-\tfrac23\kappa,\label{span}
\end{equation}
with a common constant $a$ on the domain. It is also necessary if equality
holds for every pair with nonzero $M_0$: divide the wedge equality by
$M_{0,1}M_{0,2}$ and fix one argument. Thus proportionality to $P$ is
sufficient but not necessary; a multiple of $M$ is invisible. A generic
$\bar q$ instead gives a different rank-two wedge. With nonzero $\bar m$,
the full condition is the wedge equation \eqref{T1}, not \eqref{span}.
The coefficient $-2/3$ cannot be replaced by $+1$ in this convention.

\paragraph{Darboux frame \C (\alg, \calcheck).}
With $\gamma_0=\dd\Pi_0/\dd P_0$, the corrected coefficient expressed in
physical momentum is $\gamma=\gamma_0+\hbar\partial_P\bar q$.
If $d_j$ denotes the unknown correction to the physical spatial delay,
the two Darboux shifts change by
\begin{equation}
 \delta\Delta X'_1=d_1+(\partial_{P_1}\bar q_1)M_{0,2}
                         +\gamma_{0,1}\bar m_2,\qquad
 \delta\Delta X'_2=d_2-(\partial_{P_2}\bar q_2)M_{0,1}
                         -\gamma_{0,2}\bar m_1.\label{frame}
\end{equation}
Keeping $d_j=\bar m_j=0$ arbitrarily is generally not integrable: the curl
of the prospective phase derivatives is
$-\bar q_1'M_2'-\bar q_2'M_1'$, which need not vanish.
Thus a charge correction does not by itself supply a phase generator.
Separately, conversion of an existing label-frame phase gives
\begin{equation}
 R[p]=-\sum_jp_j\partial_{P_j}\delta_0
      =\sum_jp_j\Delta X'_j=H[p]+D[p],\qquad
 D[p]=p_1\gamma_{0,1}M_{0,2}-p_2\gamma_{0,2}M_{0,1}.\label{recoil}
\end{equation}
This is not an extra term to add after already including the same argument
shift in \eqref{fixedP}. The derivative of the full classical phase, rather
than of its topological part alone, is involved.

\section{Centroid identity and the kink-height sector}
\paragraph{Local identity and boundary limit \C (\alg).}
The canonical Hamiltonian, including $\hbar\lambda\int u^2$, has
$\dot\phi=\pi$ and $u_t=\pi_x$. For a smooth compact cutoff $f_R$ define
\begin{equation}
 B_R=\int x f_R(x)u(x)\dd x,\qquad
 \dot B_R=-\int[f_R(x)+xf_R'(x)]\pi(x)\dd x.\label{cutoff}
\end{equation}
Equivalently on a finite interval,
$\partial_t\int_{-R}^R xu\dd x=[x\pi]_{-R}^R-\int_{-R}^R\pi\dd x$.
If the smeared limits exist and the last shell/boundary term vanishes,
$\dot B=-\Pi$ and $\dot M=[\pi]_{-\infty}^{\infty}=0$.
In a nonzero fixed-height sector, $M$ is a c-number and
\begin{equation}
 X_{\rm cm}=B/M,\qquad \Pi=-M\dot X_{\rm cm}.\label{centroid}
\end{equation}
There is no ordering correction to this linear continuity identity.
More generally the equal-time field algebra makes $M$ commute with $B$;
an inverse of $M$ still needs a nonzero-mass domain. Momentum eigenstates on
the line are delocalized: a localized state or a regulated matrix-element
limit is needed to define its centroid. Classical kink tails do not spoil
\eqref{cutoff}: $\pi$ is localized, although $\phi$ is not.
\D The full quantum soliton Hilbert space and these limiting domains are
not constructed by the repository's determinant packages.

\paragraph{Why dispersion is not sufficient \C/\D (\alg).}
\C If a soliton wavepacket has centroid velocity equal to its group velocity,
and if its branch energy is $E(P)=E_0(P)+\hbar F(P)$, then its first-order
charge relation is
\begin{equation}
 \bar q=-M_0\partial_PF-v_0\bar m.\label{ward}
\end{equation}
The group-velocity/state identification is \D until supplied with the
appropriate quantum sector and wavepacket. In particular, a derivative in a
fixed-$M$ sector is not automatically the derivative of a family whose
$M(P)$ changes. Even granting the branch velocity identification, omitting
$\bar m$ from \eqref{ward} is an additional assumption.

\C With fixed potential endpoints, $M=\phi(+\infty)-\phi(-\infty)$ has no
perturbative correction within that sector: decaying fluctuations integrate
to zero. This protects the endpoint definition, not the classical relation
$M_0(P_0)=2k(P_0)$. The quantum family can have a shifted relation between
height and momentum. Already the exactly solvable elastic deformation at
fixed $P$ gives
\begin{equation}
 (\bar m)_{\rm ct}=-\frac{2\lambda k}{3-4k^2},\qquad
 (\bar q)_{\rm ct}=-\frac{4\lambda kv}{3-4k^2}.\label{CTcharges}
\end{equation}
These follow by solving $P=2k^3\sqrt{1-k^2+2\hbar\lambda}/3$ to first
order, and independently from the velocity-labelled mean response below.
They demonstrate why exact topological conservation is not an argument for
$\bar m=0$ along the fixed-$P$ comparison.

\section{Independent continuum-mode insertion of the shift charge}
\paragraph{Precisely what is computed \C/\D (\modes).}
\C This section performs a new calculation of the continuum tadpole and the
linear charge insertion on a stationary, light, good soliton, with the same
free-vacuum normal ordering as the audited trace. It does not differentiate
the dispersion to obtain the answer. The final momentum result is then
compared with that derivative. \D Completeness of the nonzero-frequency
basis in the constrained quantum domain, its time-dependent projection and
the finite collective measure are not established; the result is conditional
on the scattering-continuum prescription. No claim to the fully projected
physical expectation is made.

\paragraph{Modes and local variance \C (\modes).}
Set $\xi=x-vt$, $y=\tanh(k\xi/2)$, $u_0=k^2(1-y^2)/2$,
$h_0=(1+y)/2$, $\omega=|q|\sqrt{1+q^2}$, $\nu=\omega-vq$. The potential
mode, with asymptotic normalization $1/\sqrt{2\omega}$ and measure
$\dd q/(2\pi)$, has spatial factor
\begin{align}
 f_q&=e^{iq\xi}\left[1+(T_q-1)
                  \left(h_0+\frac{h_0'}{iq}\right)\right],\\
 T_q&=-\frac{t_q-i b_q}{t_q+i b_q},\qquad
 t_q=3kq,\quad b_q=1-2k^2+2q^2-v\operatorname{sgn}(q)\sqrt{1+q^2}.
\end{align}
The co-moving pencil is checked directly at both signs of $q$ and $v$:
\begin{equation}
 [\partial_\xi^4-k^2\partial_\xi^2
    +6\partial_\xi(u_0\partial_\xi)
    +2iv\nu\partial_\xi-\nu^2]f_q=0,\qquad |T_q|=1.
\end{equation}
Write $a_q=\Re T_q-1$, $b_q^{\rm im}=\Im T_q$, and define
\begin{align}
 A_q&=\tfrac12a_q(q^2-k^2)+kq b_q^{\rm im},&
 B_q&=-a_qk^4/(8q^2),\\
 |f_q'|^2-q^2&=A_q(1-y^2)+B_qy^2(1-y^2)^2.\label{density}
\end{align}
This is a direct field-mode density, not a phase-density ansatz.
At a common infrared cutoff $\mu$, the vacuum-subtracted coincident
variance $g_\mu=\langle\eta_x^2\rangle/\hbar-\langle\eta_x^2\rangle_0/\hbar$
has the same two shapes with coefficients
\begin{equation}
 A_\mu=\int_{|q|>\mu}\frac{A_q\dd q}{4\pi\omega},\qquad
 B_\mu=\int_{|q|>\mu}\frac{B_q\dd q}{4\pi\omega}.
\end{equation}
The limit of $A_\mu$ exists. For $0<k<\sqrt3/2$, putting
$R=\sqrt{1+q^2}$ gives the independently checked integral
\begin{align}
 A&=\frac{3k^2}{4\pi}\int_1^\infty
 \left[\frac1{4R^2+4vR+4v^2-3}
       +\frac1{4R^2-4vR+4v^2-3}\right]\dd R\\
 &=\frac{\sqrt3 k}{4\pi}\operatorname{artanh}\frac{2k}{\sqrt3}.
 \label{Aintegral}
\end{align}
$B_\mu$ has a soft logarithm. It must not simply be dropped in the local
profile or in a momentum contraction.

\paragraph{Mean profile and direct $\Pi$ insertion \C (\modes).}
The order-$\hbar$ stationary field equation is
\begin{equation}
 H_v u_1=3g_\mu-2\lambda u_0,\qquad
 H_v=-\partial_\xi^2+k^2-6u_0.\label{tadpole}
\end{equation}
The following identities give an explicit solution, modulo translation:
\begin{gather}
 H_v\partial_k u_0=-2ku_0,\qquad H_vu_0''=6(u_0')^2,\\
 u_1=\left(-\frac{3A_\mu}{k^3}+\frac\lambda k\right)\partial_k u_0
                    +\frac{2B_\mu}{k^6}u_0''.\label{profile1}
\end{gather}
In particular the soft term is a translation-spreading contribution with
zero integrated height. An adjoint check uses
\begin{equation}
 w=\frac1{k^2}-\frac3{k^3}\partial_k u_0,\qquad H_vw=1,\quad
 \int w(1-y^2)\dd\xi=-\frac2{k^3},\quad
 \int wy^2(1-y^2)^2\dd\xi=0.
\end{equation}
Direct integration of \eqref{profile1}, with two distinct finite values of
$B_\mu$, verifies the same mass response. Since $\Pi=\int\pi$ is linear,
there is no quadratic Wick contraction of this operator. Its one-loop
mean correction at fixed velocity is $\int(-vu_1)$, giving
\begin{equation}
 \boxed{m_v=-\frac{3\sqrt3}{2\pi k^2}
              \operatorname{artanh}\frac{2k}{\sqrt3}+\frac{2\lambda}{k},
 \qquad q_v=-v m_v.}\label{directPi}
\end{equation}
The latter was obtained from the stationary field insertion, independently
of the energy derivative or centroid argument. Setting $q_v=0$ just because
the Gaussian contraction is zero would omit the entire mean profile.

\paragraph{Direct momentum and a nontrivial comparison \C (\modes).}
The canonical momentum has both a mean-profile term $2v\int u_0u_1$ and
a quadratic fluctuation term. Define
\begin{align}
 J_q&=\int[\Re(if_qf_q'^*)-q]\dd\xi
       =-a_q k/q+2a_q q/k+3b_q^{\rm im},\\
 G_q&=\int(|f_q'|^2-q^2)\dd\xi
       =4A_q/k+8B_q/(15k).
\end{align}
Using $\int(\partial_k u_0)(1-y^2)=2$ and
$\int(\partial_k u_0)y^2(1-y^2)^2=8/45$, the soft $B_q$ terms cancel
\emph{between} these contributions before the $q$ integral:
\begin{align}
 p_v&=\int\frac{\nu J_q-2vA_q/k}{4\pi\omega}\dd q+2\lambda vk\\
 &=-\frac{2\sqrt3 v}{\pi}\operatorname{artanh}\frac{2k}{\sqrt3}
                       +2\lambda vk.\label{directP}
\end{align}
Independent 80-digit quadratures, at five signed widths, agree with the
closed forms to better than $10^{-65}$. The mode pencil and profile equation
are separately checked. No fitted coefficient sets the tadpole result.

\C Only after this calculation, compare the audited conditional dispersion
\begin{equation}
 F_g(k;\lambda)=\frac{\sqrt3}{\pi}(\tfrac34-k^2)
 \operatorname{artanh}\frac{2k}{\sqrt3}-\frac{3k}{2\pi}
                          +\frac{2\lambda k^3}{3}.\label{Fg}
\end{equation}
It satisfies $p_v=-\partial_v F_g$. With
\begin{equation}
 V_1=-p_v/P_{0,v}=\partial_PF_g,\quad
 \bar m=m_v+M_{0,v}V_1,\quad
 \bar q=-v m_v+\Pi_{0,v}V_1,\label{directfixed}
\end{equation}
the independently computed insertion obeys \eqref{ward}, including
$-v\bar m$. At $k=.31$, $\epsilon=+1$, $\lambda=0$, for example,
\begin{equation}
 (\bar m,\bar q,V_1)=(1.2322718175732764,
                 -0.7212159462470162,-0.7263708360512509).
\end{equation}
The frozen-height prediction $-2kV_1$ even has the opposite sign here.
This compares two conditional one-body computations; it does not certify
the missing physical boundary measure.

\section{Bad/Toda calibration before the good application}
\paragraph{Continuation and dictionary \C/\D (\modes, \calcheck).}
\C Use the audited formal continuation $F_b(k)=-iF_g(ik)$ with the
action/Planck dictionary, including reversal of the finite quadratic
coefficient under continuation. For the mass coefficient the corresponding
rule is $m_b(k)=i m_g(ik)$, since both $M$ and $\hbar$ scale. Thus at
fixed velocity $v=\epsilon\sqrt{1+k^2}$,
\begin{align}
 m_{v,b}&=-\frac{3\sqrt3}{2\pi k^2}\arctan\frac{2k}{\sqrt3}
                         -\frac{2\lambda_b}{k},\\
 p_{v,b}&=-\frac{2\sqrt3 v}{\pi}\arctan\frac{2k}{\sqrt3}
                         -2\lambda_bvk,\qquad q_{v,b}=-v m_{v,b},\\
 F_b&=\frac{\sqrt3}{\pi}(\tfrac34+k^2)\arctan\frac{2k}{\sqrt3}
             -\frac{3k}{2\pi}+\frac{2\lambda_b k^3}{3},\qquad
 C_b=\lambda_b+2/\pi.\label{bad}
\end{align}
These are formal continuations, not a real whole-line bad-vacuum mode sum
\D. The lattice-matched parameter $\lambda_b=-3/(2\pi)$ gives
$C_b=1/(2\pi)$ as in the controlling one-body audit.

\C At reference spacing one, the established small-amplitude units are
\begin{equation}
 \kappa=\sqrt3\beta,\quad k=\sinh\beta,\quad
 h_T=3b\hbar_B,\quad b=e^{-1/2},\quad
 P_B=\overline P_T/(6\sqrt3b),\quad E_B=\overline E_T/(6\sqrt3b^2).
\end{equation}
The classical action obeys the corresponding scaling, so its dimensionless
phase is unchanged. Put $k_1=\rho$, $k_2=r\rho$, $0<r<1$, on the
right-moving overtaking branch. The same $r$ is the leading
$\sqrt{y_2/y_1}$ or cube root of the physical momentum ratio.
\D The additional identification tested here is between the
continuum canonical \emph{recoil representative} and the exchange-odd part
of Toda's two-term \emph{directional residual}, using their stated classical
subtractions. There is no established common quantum operator/endpoint
dictionary identifying those entire observables.

\paragraph{What actually reproduces the rational term \C (\calcheck).}
With $\sigma=+1$ for good and $\sigma=-1$ for bad, put
$C=\lambda-2\sigma/\pi$. At small amplitude the canonical insertion and
Darboux coefficient obey
\begin{equation}
 p_v=2\sigma Ck+O(k^3),\qquad \gamma_0=-k^{-2}+O(1),\qquad
 D[p]\longrightarrow-4\sigma C(r-r^{-1}).\label{Dlimit}
\end{equation}
Consequently the matched bad recoil is exactly $2(r-r^{-1})/\pi$.
Independent regular quadrature verifies the audited decomposition
\begin{align}
 R_{12}^{(1)}&=\frac4\pi\int_0^r\log(s/r)\frac{s^2}{1-s^2}\dd s
             =-\frac4\pi[\chi_2(r)-r],\\
 S^{(1)}&=\frac2\pi(r+r^{-1})-\pi,\\
 \vartheta_N^{(1)}&=R_{12}^{(1)}-S^{(1)}
       =\pi-\frac4\pi\chi_2(r)+\frac2\pi(r-r^{-1}),\qquad
 \chi_2(r)=\tfrac12[\operatorname{Li}_2(r)-\operatorname{Li}_2(-r)].
 \label{Toda}
\end{align}
At $r=1/2$, the rational part is $-0.9549296585513720$.
The whole $\vartheta_N$ is $1.5305278132789631$. Recoil accounts for neither
the chi-function term nor the endpoint constant.
\C As an additional check against overidentification, the audited canonical
continuum candidate $J+D$ differs from \eqref{Toda} by
\begin{equation}
 (J+D)-\vartheta_N=-\frac\pi2-\frac2\pi\log r\log\frac{1-r}{1+r},
\end{equation}
a nonconstant gap, $-2.0555822900856296$ at $r=1/2$. Rational agreement
is not validation of the full phase.

\paragraph{Why this is not $c_{\rm bad}$ \C (\calcheck, \modes).}
The proposed finite-coefficient residual scales as
\begin{equation}
 Q_{12}=\tfrac43r(1-r^2)\rho^4+O(\rho^6).
\end{equation}
It cannot reproduce an order-one rational function of $r$. Conversely,
every finite $c$ disappears from this leading Toda limit, so that limit
cannot determine it. Dividing the rational term by $Q_{12}$ produces a
kinematic divergence, not a coefficient in the assumed loop expansion.

\C The actual conditional charge response makes the problem still more
explicit. For either sign of the equation, for $\epsilon=+1$,
\begin{align}
 \bar m&=\frac1{\pi k}+\left(\frac{16\sigma}{9\pi}-\frac{2\lambda}{3}\right)k
           +O(k^3),\\
 \bar q&=-\frac1{\pi k}+\left(\frac{49\sigma}{18\pi}-\frac{4\lambda}{3}\right)k
           +O(k^3).\label{smallcharges}
\end{align}
Changing $\Pi$ alone in $T_0$ would yield $4(r-r^{-1})/(3\pi)$, only
two thirds of Toda's rational coefficient. The simultaneous $M$ term in
\eqref{T1} yields its negative. The surviving topological substitution is
\begin{equation}
 T_1|_P=-\frac{2\sigma}{3\pi}\rho^2(1-r^2)(r+r^{-1})+O(\rho^4),
 \label{toplimit}
\end{equation}
not a constant times $Q_{12}$ and not \eqref{Dlimit}. The leading
$1/k$ charge shift is incompatible with $\Delta\Pi=\hbar\kappa P$.
These are sequential one-loop coefficients; their use requires
$\hbar/k^2\ll1$. No uniform finite-$\hbar$ zero-amplitude limit is claimed.

\section{Good equation: parameter dependence and the failed frozen-height shortcut}
\paragraph{What can be given as a function of $\lambda$ \C (all but \proj).}
Equations \eqref{directPi}, \eqref{directP} and \eqref{directfixed} give the
conditional one-body response for every finite $\lambda$. The audited
small-amplitude coefficient and the isolated rational recoil are
\begin{equation}
 C_g(\lambda)=\lambda-2/\pi,\qquad
 D_g[p]\longrightarrow(-4\lambda+8/\pi)(r-r^{-1}).\label{goodrecoil}
\end{equation}
\begin{center}
\begin{tabular}{@{}cccc@{}}\toprule
$\lambda$ & $C_g$ & Coefficient in \eqref{goodrecoil} & Physical $c_g$\\\midrule
$9/(4\pi)$ & $1/(4\pi)$ & $-1/\pi$ & undetermined\\
$0$ & $-2/\pi$ & $8/\pi$ & undetermined\\\bottomrule
\end{tabular}
\end{center}
\D These recoil coefficients are not values of $c_g$. They neither select
a scheme nor complete the physical phase.

\C Counterterm accounting is an independent safeguard against double
counting. In the same specified chamber frame, the complete classical
elastic response at fixed momenta is
\begin{equation}
 \delta_1^{\rm ct}=\lambda K_\lambda,\qquad
 K_\lambda=3\delta_0-4\sum_jP_j\partial_{P_j}\delta_0.
\end{equation}
Fresh finite differences of the exact stiffness-deformed phase, with both
momenta inverted anew, verify $K_\lambda^{++}=0.05008140486841891$ and
$K_\lambda^{+-}=-0.22517393868334176$ at $\alpha=(.3,.6)$.
The isolated $\lambda$ dependence of \eqref{goodrecoil} must be accompanied
by the corresponding change in the label-frame contribution; replacing
the momentum convention while leaving that contribution fixed is invalid.
No $\lambda K_\lambda$ already present in a representative may be added again
under the name of a new CDD coefficient.

\paragraph{Even imposing $\bar m=0$ does not rescue a universal $c$ \C (\alg).}
If one nevertheless grants both the branch-velocity identification and
the frozen-height assumption, then
\begin{equation}
 \bar q_{\rm fr}=-M_0 F_P,\qquad
 c_{\rm eff}(P_1,P_2)=\frac23\frac{F_P(P_1)-F_P(P_2)}{P_1/M_1-P_2/M_2}.
\end{equation}
This is a pointwise quotient, not a constant. For $\lambda=0$ on the good
branch, it is $-2.4463194544$ at $(k_1,k_2)=(.2,.4)$ overtaking,
$-4.8131022786$ at $(.3,.6)$ overtaking, and $-15.8575027361$ at
$(.2,.4)$ head-on. The corresponding three values at
$\lambda=9/(4\pi)$ are $-0.7124456976$, $-1.5929611493$, and
$1.1003496931$. The failure is structural, not a numerical tolerance issue.

\C On a single right-moving branch only, its small-width limit exists.
Expanding $F_P=C+D_2k^2+O(k^4)$ gives
\begin{equation}
 D_2=\frac{5\sigma\lambda}{6}-\frac{23}{9\pi},\qquad
 c_{{\rm fr},\rm small}=2D_2.
\end{equation}
The resulting good values are $-49/(36\pi)$ at $9/(4\pi)$ and
$-46/(9\pi)$ at zero; the matched bad value is $-47/(18\pi)$.
They characterize this expressly false-for-the-computed-branch
frozen-height shortcut, not $c_{\rm good}$ or $c_{\rm bad}$.
In opposite directions even the common leading velocity correction changes
sign, so it cannot be discarded as one common multiple of $M$ in
\eqref{span}.

\section{Ring and one collision-point projection checks}
\paragraph{Conditional ring correction \C (\calcheck).}
For the proposed directional ring phase $\Phi(j\mid l)=\delta(j,l)+\Pi_jM_l$,
the twist changes by
\begin{equation}
 (\Pi_jM_l)_1|_P=\bar q_jM_{0,l}+\Pi_{0,j}\bar m_l.
\end{equation}
The symmetric pair correction is exactly \eqref{S1}, and summing the
assumed equations, including the audited self-winding term, changes the
total twist by
\begin{equation}
 (\Pi_{\rm tot}M_{\rm tot})_1
       =\bar q_{\rm tot}M_{0,\rm tot}+\Pi_{0,\rm tot}\bar m_{\rm tot}.
\end{equation}
The same charge/height response therefore appears in this algebraic check.
An antisymmetric $cQ_{jl}$ cancels from the total pair sum. \D The ring
holonomy is still not a quantum wavefunction-transport calculation; it
provides no independent equation selecting $c$.

\paragraph{Exact two-soliton background at one point \C (\proj).}
At $\alpha_1=.3$, $\alpha_2=.6$, $\epsilon_1=\epsilon_2=+1$,
$t=X_1=X_2=0$, the checker uses the exact two-soliton tau function,
$\tau=1+E_1+E_2+A_{12}E_1E_2$, $E_j=e^{k_j(x-X_j)-k_jv_jt}$,
and differentiates $\phi=2\partial_x\log\tau$, $\pi=\partial_t\phi$.
It evaluates the four-modulus symplectic Gram matrix directly, retaining
all cross terms. Here
\begin{equation}
 A_{12}=0.09096862155680799,\qquad
 \det W=(P_{0,\alpha_1}P_{0,\alpha_2})^2
          =0.0032005493687711533.
\end{equation}
The determinant error is below $10^{-50}$ at 65 digits. For the localized
test fluctuation $\eta=(0,e^{-x^2})$, its symplectic projection onto these
four tangents has
\begin{equation}
 (\mathcal P\eta)_\phi(-\infty)=0,\qquad
 (\mathcal P\eta)_\phi(+\infty)=-1.0756732271416573.
\end{equation}
Thus the complementary radiation projection does not preserve fixed
potential endpoints even at this particular collision point.

\D This is an instantaneous Gram determinant, explicitly \emph{not} the
requested spacetime projected fluctuation determinant. A numerical value
of the latter would require a boundary-compatible evolving domain, the
constraint/collective-radiation Schur complement, and time-endpoint state
normalization. Those inputs are not supplied by the audited machinery.
Rather than choose them by preference or assign the determinant a zero
phase, the cache records \file{null}. A single kinematic point does not
remove this missing-definition problem.

\section{Claim summary, failures recorded, and conclusion}
\small
\begin{longtable}{@{}L{5.3cm}L{6.1cm}L{3.4cm}@{}}\toprule
Claim and status & Evidence and limit & Checker\\\midrule\endhead
\C One coefficient in the exhibited $P,M$ wedge; defect rank zero & Polynomial primitives, mixed minors and explicit transmission lift. Not a classification. & \alg\\[3pt]
\C Fixed-$P$ charge and chamber algebra & Symbolic inversion, symmetric subtraction, curl counterexample, exact elastic deformation & \alg\\[3pt]
\C Continuity/centroid implication & Regulated integration by parts; quantum domain and group-velocity identification remain conditional & \alg\\[3pt]
\C Conditional direct $m_v,q_v,p_v$ & Field-mode densities, pencil test, explicit mean-profile solution, independent 80-digit integrals & \modes\\[3pt]
\C Bad/Toda rational recoil match & Symbolic small-width limit and independent ratio-integral quadrature; full phase and $c$ not matched & \calcheck\\[3pt]
\C Good $\lambda$ dependence & Conditional one-body responses and isolated recoil; full counterterm response checked separately & \alg, \modes, \calcheck\\[3pt]
\C Conditional ring sum and collision geometry & Charge-product expansion; direct exact-tau symplectic integration and endpoint failure & \calcheck, \proj\\[3pt]
\D Physical residual $c$, projected expectation and functional determinant & Missing domain/state map and dynamical two-body calculation; explicit nulls & Open\\\bottomrule
\end{longtable}
\normalsize

\paragraph{Numerical scope \C.}
The scripts use 80 decimal digits except the 65-digit collision quadrature.
Analytic/integral comparisons are asserted at $10^{-65}$ or $10^{-70}$;
the finite-width Toda recoil test uses $k_1=10^{-16}$ and a distinct
$10^{-29}$ truncation tolerance. Elastic-response Richardson differences
use steps $10^{-14}$ and half that, with $10^{-50}$ tolerance.
The collision integrates to $|x|=600$ and separately verifies its endpoint
limit. These are numerical diagnostics, not rigorous interval error bounds.
All four scripts print \texttt{ALL CHECKS PASSED}; their cached runtimes
total approximately 35 seconds, far below thirty minutes. Physical interpretation
errors are not hidden inside those numerical tolerances.

\paragraph{Failures recorded.}
\begin{enumerate}
\item \C The initial stronger premise that the full residual has only one
coefficient is unsupported: the defect package exhibits one wedge while
explicitly retaining other unknown contributions. The work was adapted to
that actual scope.
\item \C The frozen-height dispersion shortcut fails constant-$c$ and
cross-channel tests; the nonzero direct height response is retained.
\item \C The Gaussian linear-$\Pi$ contraction alone misses the mean-profile
term. The unprojected local tadpole has a soft logarithm. Its integrated
mass contribution vanishes and its momentum contribution cancels by exact
identities; no interval or nonzero contribution was discarded.
\item \C Changing the shift charge alone gives a rational coefficient
of the wrong size for Toda and omits an equal, opposite height term.
The full topological substitution and $P\wedge M$ have different scaling.
\item \C A Darboux-coefficient change with unchanged spatial delays fails
the curl condition. The known rational recoil cannot be counted a second
time as a newly computed residual.
\item \C The conditional continuum/Toda full-phase comparison retains
a nonconstant gap after the rational terms agree.
\item \C At the single collision point the instantaneous projection fails
the fixed-endpoint test. A physical two-body functional determinant and
quantum ring holonomy were not evaluated, and are not reported as zero.
\item \C No implementation assertion failed in these four scripts, and
no tolerance was weakened. The first PDF build failed because the checker
macro name conflicted with the standard LaTeX command \file{cal}; it was
renamed before rebuilding. A second build found an accidental form-feed
character in a formula command; the character was repaired and the long
charge display was split for layout. The following build exposed two
missing backslashes before paired left delimiters; those were restored. The two optional audit-file reads reported
missing files; no claim about their unobserved verdicts is made.
\end{enumerate}

\paragraph{Conclusion (C).}
\C The continuity identity survives the stated boundary/domain assumptions;
the conditional direct mode insertion agrees with it only after the height
response is included. The tested mechanism does not isolate a universal
$c(P_1M_2-P_2M_1)$. Toda calibrates the already-known rational recoil,
not that coefficient. \D Therefore $c_{\rm bad}$ and
$c_{\rm good}(\lambda)$, including its values at $0$ and $9/(4\pi)$,
remain undetermined. What is missing is precisely a common physical
endpoint/state map, the boundary-compatible collective projection and
two-body determinant (or equivalent dynamical phase input). No additional
bootstrap identity, choice of the two candidate $\lambda$ values, or
centroid Ward identity supplied here replaces those inputs.

\paragraph{Internal references and reading scope.}
{\raggedright
The following repository texts, rather than newly verified literature
entries, supply the inherited targets and qualifications:
\file{README.md}, \file{AGENTS.md}, \file{report.tex}, \file{errata.tex};
\file{good_twobody.tex}, \file{audit_good_twobody.tex} (frame and constraint
passages); \file{defect_bootstrap.tex}, \file{audit_resonant_channels.tex};
\file{good_renorm.tex}, \file{onebody_scheme.tex},
\file{audit_onebody_scheme.tex}, \file{onebody_dispersion_test.tex};
\file{det_projected_oneloop.tex}, \file{audit_det_projected.tex},
\file{charge_frame_master.tex}, \file{audit_charges_action.tex};
\file{toda_route.tex}, \file{twobody_toda.tex},
\file{audit_twobody_toda.tex}, \file{audit_frame_dictionary.tex},
\file{audit_twobody_cont.tex}; \file{smx_candidates.tex}, \file{audit_smx.tex}.
Relevant existing script implementations were read only; all reported new
computations use the four independent \file{shiftcharge} scripts.\par}
\end{document}
